\chapter{Marco referencial}

En este marco, usted debe analizar todo lo que se ha escrito, desarrollado e
investigado sobre el objeto de su estudio.
Tenga en cuenta: lo que se sabe del tema, los estudios que se han realizado,
desde qué puntos se han tratado. Cuando se desarrolla el marco referencial o
teórico, se sustenta de forma teórica el problema de su proyecto, para ello, no
se limite únicamente a describirlo, sino que deberá revisar diferentes teorías y
aspectos de la misma. Procure que el objeto de su estudio sea esclarecido por
medio de las teorías que aborde, explique cómo y cuándo ocurre, establezca
relaciones internas y externas del fenómeno que investiga y defina de forma
clara cómo enfocará el problema para su estudio.

\section{Marco teórico}

En el marco teórico se reúne información documental con el fin de establecer
el cómo y qué información se acopie, de qué manera se analizará y estimar el
tiempo en que se demore el desarrollo.
La información seleccionada en el marco teórico permite tener un
conocimiento profundo del tema en estudio y que le da significado a la
investigación y/o desarrollo, de aquí se pueden generar nuevos conocimientos.
Recuerde, de las teorías que apoyen su proyecto dependerá la validez interna
y externa.

\section{Antecedentes o estado del arte}

Se revisan publicaciones recientes para ver si se encuentra algún proyecto
hecho en otra parte o en otro momento que se parezca a la que se esta
desarrollando y así examinar sus resultados y forma de enfocarla, de manera
que no se cometan los mismos errores, o simplemente aprovechar de ellos lo
que le sirva y oriente en el desarrollo del proyecto.

\section{Marco legal}
\lipsum[1]

\begin{table}[tb]
    \caption{Tabla ejemplo}
    \label{tab:ejemplo}
    \centering
    \begin{tabular}{ccc}
        \toprule
        nombre & peso  & comida  \\
        \midrule
        ratón  & 10 g  & queso   \\
        gato   & 1 kg  & ratones \\
        perro  & 10 kg & gatos   \\
        t-rex  & 10 Mg & perros  \\
        \bottomrule
    \end{tabular}
\end{table}

Puede hacer referencia a la tabla \ref{tab:ejemplo}, a las figuras \ref{fig:ipsum} y a la figura \ref{fig:uan}.

\begin{figure}[tb]
    \centering
    \includegraphics[width=0.5\columnwidth]{uan.png}
    \caption[A floating figure]{A floating figure (the lithograph \emph{Galleria di stampe}, of M.~Escher, got from \url{http://www.mcescher.com/}).}
    \label{fig:uan}
\end{figure}

\begin{figure}[tb]
    \centering

    \subfloat[Asia personas duo.]
    {\includegraphics[width=.45\columnwidth]{uan.png}} \quad
    \subfloat[Pan ma signo.]
    {\label{fig:ipsum}%
        \includegraphics[width=.45\columnwidth]{uan.png}} \\
    \subfloat[Methodicamente o uno.]
    {\includegraphics[width=.45\columnwidth]{uan.png}} \quad
    \subfloat[Titulo debitas.]
    {\includegraphics[width=.45\columnwidth]{uan}}
    \caption[Tu duo titulo debitas latente]{Tu duo titulo debitas
        latente.}
    \label{fig:esempio}
\end{figure}

